% \begin{center}
% \textbf{\large{Abstract}}
% \end{center}

% \addcontentsline{toc}{chapter}{Abstract}
% \begin{small}

% This report presents the analysis, design, and implementation of \textbf{Fitty}, an Android mobile application designed to support gym-oriented workout planning and personalized health tracking. The core idea of the system is to act as a digital fitness companion that not only lists exercises, but also helps users learn proper movements, build and maintain a suitable training schedule, monitor body changes, and record nutrition-related data during their training journey.

% To address this problem, the system is implemented using a layered architecture. Kotlin and Jetpack Compose are used for the presentation layer, a dedicated business layer coordinates the main flows, Room and DataStore are employed for local persistence, and Firebase provides cloud-based services for authentication, data storage, media storage, and push notifications. A notable technical aspect of the system is the adoption of an \textit{offline-first} approach for the exercise module, where exercise metadata is synchronized from Firestore, validated, normalized, and stored locally so that the user interface can access it more efficiently and reliably. In addition, the application supports dynamic content management, allowing onboarding content, exercise categories, starter plan templates, and behavioral configurations to be updated without modifying the user interface source code.

% The outcome of the report is a functional prototype of a personalized fitness support application for gym users. The implemented system includes the main modules for sign-up and sign-in, guest access, initial profile setup and plan preview, workout scheduling, exercise browsing by muscle category, exercise detail and instructional media viewing, quick workouts and scheduled workouts, body and nutrition tracking, user profile management, local reminders, push notifications, and AI-assisted meal and body-scan analysis. The project also provides unit tests for several important business components such as starter plan generation, exercise synchronization, workout completion, and progress tracking.

% Overall, the report demonstrates that combining a modern Android architecture with cloud services, local persistence, and supporting AI services is a suitable approach for building personalized fitness applications. The achieved results provide a solid foundation for future extensions such as deeper nutrition planning, richer body-progress analytics, improved coaching logic, and a more complete real-world training experience.

% \vspace*{1cm}
% \textbf{Keywords}: Android, Jetpack Compose, Firebase, offline-first, personalized fitness application
% \end{small}
